\chapter{Τεχνική Υλοποίηση Συστήματος}
Το παρόν κεφάλαιο αναλύει την τεχνική υλοποίηση του ψηφιακού διδύμου, παρουσιάζοντας το σύστημα όχι ως στατική κατασκευή, αλλά ως αποτέλεσμα μιας εξελικτικής πορείας δύο διακριτών φάσεων. Η πρώτη φάση αντιστοιχεί στο τοπικό πρωτότυπο (\textlatin{Proof of Concept} ή \textlatin{MVP}), μέσω του οποίου επιβεβαιώθηκε η αξιοπιστία της συλλογής δεδομένων, η λειτουργία των βασικών αυτοματισμών και η δυνατότητα \textlatin{closed-loop} ελέγχου. Η δεύτερη φάση αντιστοιχεί στην ωρίμανση του συστήματος προς μια πλήρως προσαρμοσμένη πλατφόρμα \textlatin{Enterprise Cloud SaaS}, η οποία μεταφέρει το επίκεντρο από το \textlatin{engineering console} περιβάλλον σε ένα \textlatin{user-centric}, \textlatin{cloud-native} και \textlatin{responsive} οικοσύστημα.

Η δομή του κεφαλαίου ακολουθεί αυτή ακριβώς τη λογική. Αρχικά περιγράφεται η τοπική υλοποίηση που στηρίχθηκε σε \textlatin{MQTT}, \textlatin{openHAB}, \textlatin{InfluxDB} και \textlatin{Grafana}. Στη συνέχεια παρουσιάζεται η νέα \textlatin{cloud-based} αρχιτεκτονική, με \textlatin{frontend} σε \textlatin{React 18} και \textlatin{backend} σε \textlatin{Supabase/PostgreSQL}, καθώς και η τρισδιάστατη οπτικοποίηση και οι λειτουργίες διαλειτουργικότητας και αλγοριθμικής επέκτασης που ολοκλήρωσαν την τελική μορφή της πλατφόρμας.

\section{\texorpdfstring{Φάση 1: Τοπική Υλοποίηση και Συλλογή Δεδομένων (\textlatin{MVP})}{Φάση 1: Τοπική Υλοποίηση και Συλλογή Δεδομένων (MVP)}}
Η πρώτη υλοποίηση της πλατφόρμας αναπτύχθηκε με στόχο να αποδείξει ότι ένα σχολικό περιβάλλον μπορεί να αποκτήσει λειτουργίες Ψηφιακού Διδύμου μέσα από χαμηλού κόστους \textlatin{edge} συσκευές, \textlatin{semantic middleware} και βασικούς μηχανισμούς οπτικοποίησης. Σε αυτό το στάδιο το σύστημα ήταν τοπικά αναπτυγμένο, με τον \textlatin{Raspberry Pi} να λειτουργεί ως κεντρικός κόμβος διαμεσολάβησης, συλλογής τηλεμετρίας και εκτέλεσης κανόνων.

Το τοπικό πρωτότυπο ήταν εξαιρετικά σημαντικό για δύο λόγους. Πρώτον, επέτρεψε τη σταδιακή βαθμονόμηση του \textlatin{edge hardware}, των \textlatin{topics}, των κανόνων και των \textlatin{thresholds} σε πραγματικές συνθήκες πεδίου. Δεύτερον, παρήγαγε την πρώτη αξιόπιστη μάζα ιστορικών δεδομένων πάνω στην οποία στηρίχθηκαν τόσο η ενεργειακή ανάλυση όσο και η μετέπειτα σχεδίαση της τελικής \textlatin{cloud} πλατφόρμας. Με άλλα λόγια, η αρχική υλοποίηση δεν αποτελεί «ξεπερασμένο» στάδιο, αλλά τον αποδεικτικό πυρήνα της λειτουργικής αξίας του συστήματος.

\subsection{\texorpdfstring{Τοπική Δικτυακή Υποδομή με \textlatin{MQTT} και \textlatin{Eclipse Mosquitto}}{Τοπική Δικτυακή Υποδομή με MQTT και Eclipse Mosquitto}}
Η επικοινωνία μεταξύ των \textlatin{edge} συσκευών και του κεντρικού εξυπηρετητή υλοποιήθηκε αρχικά αποκλειστικά μέσω \textlatin{MQTT}, με τοπική εγκατάσταση του \textlatin{Eclipse Mosquitto} σε \textlatin{Raspberry Pi}, ο οποίος ανέλαβε τη δρομολόγηση όλων των ροών τηλεμετρίας και ελέγχου \cite{EclipseMosquitto2024}. Στο επίπεδο παραμετροποίησης, κάθε κόμβος δημοσίευε τις μετρήσεις του σε προκαθορισμένα \textlatin{topics}, ενώ το \textlatin{openHAB} και τα υπόλοιπα υποσυστήματα εγγράφονταν στα αντίστοιχα κανάλια για λήψη μετρήσεων ή αποστολή εντολών.

Η δομή των \textlatin{MQTT topics} οργανώθηκε ιεραρχικά ώστε να αντιστοιχεί στη φυσική εγκατάσταση και στα επιμέρους υποσυστήματα. Έτσι, οι ροές θορύβου, ισχύος, θερμοκρασίας, κατάστασης θυρών και εντολών ενεργοποίησης αποτυπώθηκαν σε διακριτά \textlatin{topics} ανά χώρο και συσκευή, ώστε να υποστηρίζεται καθαρή δρομολόγηση δεδομένων και εύκολη επέκταση του σχήματος ονοματοδοσίας.

\subsection{\texorpdfstring{\textlatin{Semantic Middleware} και Κανόνες Ελέγχου στο \textlatin{openHAB}}{Semantic Middleware και Κανόνες Ελέγχου στο openHAB}}
Στο δεύτερο επίπεδο του πρωτοτύπου χρησιμοποιήθηκε το \textlatin{openHAB} ως \textlatin{semantic middleware} και \textlatin{rules engine}. Στο πλαίσιο της πρακτικής υλοποίησης, το \textlatin{openHAB} συγκέντρωνε τη ροή δεδομένων από \textlatin{MQTT}, τη μετέτρεπε σε δομημένα \textlatin{Items}, αντιστοιχισμένα σε συγκεκριμένους χώρους και συσκευές, και εφάρμοζε βασικούς αυτοματισμούς όπως:
\begin{itemize}
  \item τον έλεγχο θορύβου στην αίθουσα μέσω της \textlatin{quiet-lamp},
  \item την ανίχνευση ανοιχτής πόρτας ψυγείου και την ειδοποίηση για λειτουργική ανωμαλία,
  \item την αρχική εποπτεία ενεργειακών φορτίων στο κυλικείο.
\end{itemize}

Η παραμετροποίηση αυτή υλοποιήθηκε με κανόνες που χρησιμοποιούσαν επιχειρησιακά \textlatin{thresholds}, χρονόμετρα και συνθήκες κατάστασης. Ενδεικτικά, ο κανόνας \textlatin{NoiseGuard} ενεργοποιούσε την \textlatin{quiet-lamp} όταν η στάθμη θορύβου υπερέβαινε το καθορισμένο όριο, ενώ ο \textlatin{FridgeGuardian} συνέδεε την κατάσταση της θύρας με χρονική καθυστέρηση και ενεργοποίηση \textlatin{buzzer} ειδοποίησης (βλ.\ Παράρτημα Γ, Λίστα Γ.1 και Γ.2 για τον αναλυτικό κώδικα των κανόνων \textlatin{NoiseGuard} και \textlatin{FridgeGuardian} αντίστοιχα).

\subsection{\texorpdfstring{Ιστορική Καταγραφή σε \textlatin{InfluxDB} και Αρχικά \textlatin{Dashboards} σε \textlatin{Grafana}}{Ιστορική Καταγραφή σε InfluxDB και Αρχικά Dashboards σε Grafana}}
Για την αποθήκευση χρονοσειρών και την ιστορική τεκμηρίωση της λειτουργίας του συστήματος, το τοπικό πρωτότυπο διασυνδέθηκε με \textlatin{InfluxDB} \cite{InfluxData2024}. Η βάση αυτή επιλέχθηκε επειδή ήταν ιδιαίτερα κατάλληλη για ταχεία εγγραφή τηλεμετρίας και ανάλυση μετρήσεων που μεταβάλλονται χρονικά, όπως ο θόρυβος, η ενεργός ισχύς και οι καταστάσεις εξοπλισμού.

Πάνω από αυτό το επίπεδο αναπτύχθηκαν αρχικά \textlatin{dashboards} σε \textlatin{Grafana}, τα οποία παρείχαν γραφήματα χρονοσειρών, οπτική επιβεβαίωση ανωμαλιών και εποπτεία της ενεργειακής συμπεριφοράς του συστήματος \cite{GrafanaLabs2024}. Τα \textlatin{dashboards} αυτά υπήρξαν καθοριστικά για την πρώτη ανάλυση των δεδομένων και την ποσοτική τεκμηρίωση των \textlatin{KPIs} της μελέτης. Ειδικότερα, μέσα από αυτά κατέστη εφικτή η οπτική αναγνώριση υπερβάσεων θορύβου, η σύγκριση φορτίων πριν και μετά από παρεμβάσεις και η δημιουργία ιστορικών αναφορών.

\subsection{Αποδεικτική Αξία και Περιορισμοί του Τοπικού Πρωτοτύπου}
Η πρώτη αυτή φάση απέδειξε ότι η συλλογή δεδομένων, η σημασιολογική μοντελοποίηση και οι βασικοί αυτοματισμοί μπορούσαν να λειτουργήσουν αξιόπιστα σε πραγματικό σχολικό περιβάλλον. Ωστόσο, ανέδειξε και σαφείς περιορισμούς. Οι διαθέσιμες διεπαφές ήταν κυρίως \textlatin{engineering-oriented} και όχι σχεδιασμένες για μη τεχνικούς χρήστες. Τα \textlatin{dashboards} και τα \textlatin{floorplan views} ήταν επαρκή για τεχνική εποπτεία, αλλά όχι ιδανικά για επιχειρησιακή χρήση από διευθυντές ή διοικητικό προσωπικό. Επιπλέον, η απομακρυσμένη πρόσβαση στο πρωτότυπο περιβάλλον ήταν περισσότερο σύνθετη και προϋπέθετε μηχανισμούς ασφαλούς τεχνικής σύνδεσης, γεγονός που περιόριζε τη φυσικότητα της εμπειρίας για τον τελικό χρήστη.

Με άλλα λόγια, η Φάση 1 απέδειξε την αξία της συλλογής, της σημασιολογικής οργάνωσης και του \textlatin{closed-loop automation}, αλλά κατέστησε εμφανές ότι το επόμενο βήμα έπρεπε να είναι η ανάπτυξη ενός νέου επιπέδου υπηρεσιών, ειδικά σχεδιασμένου για \textlatin{cloud} πρόσβαση, κλιμάκωση και βελτιωμένο \textlatin{UI/UX}. Η λογική κατάσταση του υποσυστήματος ελέγχου ψυγείου, όπως υλοποιήθηκε στο αρχικό \textlatin{proof-of-concept} επίπεδο, παρουσιάζεται στο Σχήμα~\ref{fig:fridge_fsm}, ενώ στο Σχήμα~\ref{fig:openhab_dashboard_legacy} αποτυπώνεται η αρχική \textlatin{engineering-oriented} διεπαφή του \textlatin{openHAB}. Το Σχήμα 4.1 απεικονίζει τη Μηχανή Πεπερασμένων Καταστάσεων (\textlatin{Finite State Machine}) του κανόνα \textlatin{FridgeGuardian}. Το σύστημα ξεκινά από την κατάσταση \textlatin{S0} (Πόρτα Κλειστή -- κανονική λειτουργία). Μόλις ο αισθητήρας \textlatin{ESP8266} ανιχνεύσει άνοιγμα της θύρας, μεταβαίνει στην κατάσταση \textlatin{S1} (Πόρτα Ανοιχτή) όπου ενεργοποιείται χρονόμετρο. Εάν η πόρτα παραμείνει ανοιχτή για περισσότερο από 45 δευτερόλεπτα, το σύστημα μεταβαίνει στην κατάσταση \textlatin{S2} όπου ενεργοποιείται το \textlatin{buzzer} ειδοποίησης. Το κλείσιμο της πόρτας οποιαδήποτε στιγμή επαναφέρει το σύστημα στην κατάσταση \textlatin{S0}.

\begin{figure}[htbp]
\centering
\includegraphics[width=\textwidth]{fig_4_1}
\caption{Απεικόνιση της λογικής κατάστασης του υποσυστήματος ελέγχου ψυγείου, όπως υλοποιήθηκε στο αρχικό \textlatin{proof-of-concept} επίπεδο.}
\label{fig:fridge_fsm}
\end{figure}

\begin{figure}[htbp]
\centering
\includegraphics[width=\textwidth]{openhab_dashboard}
\caption{Ενδεικτική απεικόνιση του αρχικού \textlatin{engineering-oriented dashboard} του \textlatin{openHAB}, όπως χρησιμοποιήθηκε στο στάδιο του τοπικού πρωτοτύπου πριν από τη μετάβαση στη τελική \textlatin{cloud} διεπαφή.}
\label{fig:openhab_dashboard_legacy}
\end{figure}

\section{\texorpdfstring{Φάση 2: Ανάπτυξη \textlatin{Enterprise Cloud Platform (SaaS)}}{Φάση 2: Ανάπτυξη Enterprise Cloud Platform (SaaS)}}
Η δεύτερη φάση ανάπτυξης συνιστά τη μετάβαση από το τοπικό πρωτότυπο σε μια ώριμη \textlatin{cloud-based} πλατφόρμα, σχεδιασμένη να λειτουργεί ως ενιαίο περιβάλλον υπηρεσιών \textlatin{Smart Campus}. Σε αυτή τη φάση το βάρος μεταφέρθηκε από την τεχνική επαλήθευση του \textlatin{middleware} προς τη συνολική εμπειρία χρήστη, τη διαθεσιμότητα 24/7, τη δυνατότητα πρόσβασης από πολλαπλές συσκευές και την παροχή περισσότερο συγκεντρωτικών και επιχειρησιακά αξιοποιήσιμων προβολών.

Η νέα αρχιτεκτονική δεν ακυρώνει το τοπικό επίπεδο. Αντιθέτως, το διατηρεί ως αξιόπιστο μηχανισμό επικοινωνίας με το \textlatin{hardware}, ενώ μεταθέτει το επίπεδο παρουσίασης, υπηρεσιών και διαχείρισης σε \textlatin{cloud} περιβάλλον. Το \textlatin{openHAB} παραμένει συνεπώς κομβικό ως \textlatin{semantic middleware}, αλλά δεν αποτελεί πλέον το τελικό \textlatin{UI} της πλατφόρμας.

\subsection{\texorpdfstring{\textlatin{Frontend} με \textlatin{React 18, Vite} και \textlatin{Tailwind CSS}}{Frontend με React 18, Vite και Tailwind CSS}}
Το \textlatin{frontend} της νέας πλατφόρμας υλοποιήθηκε ως σύγχρονη \textlatin{Single-Page Application} βασισμένη σε \textlatin{React 18}, με εργαλείο δόμησης το \textlatin{Vite} και με χρήση \textlatin{Tailwind CSS} για τον οπτικό και \textlatin{responsive} σχεδιασμό της διεπαφής. Η επιλογή αυτής της στοίβας υπαγορεύθηκε από την ανάγκη για γρήγορη απόδοση στον \textlatin{browser}, αρθρωτή σύνθεση της διεπαφής και δυνατότητα ανάπτυξης εξειδικευμένων προβολών ανά ρόλο, υποσύστημα και σενάριο χρήσης.

Η νέα διεπαφή δεν αναπαράγει τη λογική ενός εργαλείου αυτοματισμού. Σχεδιάστηκε ώστε να υποστηρίζει λειτουργίες διοικητικής εποπτείας, παρακολούθησης κρίσιμων δεικτών, διαχείρισης συσκευών και πρόσβασης σε \textlatin{2D/3D} προβολές του σχολικού συγκροτήματος μέσα από ένα συνεκτικό περιβάλλον. Αυτή η μετατόπιση από \textlatin{system-centric} σε \textlatin{user-centric UX} αποτελεί έναν από τους βασικούς μετασχηματισμούς της πλατφόρμας.

\subsection{\texorpdfstring{\textlatin{Backend} με \textlatin{Supabase, PostgreSQL} και \textlatin{Edge Functions}}{Backend με Supabase, PostgreSQL και Edge Functions}}
Στο \textlatin{backend} επίπεδο η πλατφόρμα μεταφέρθηκε σε υποδομή νέφους με \textlatin{Supabase} και \textlatin{PostgreSQL}. Η επιλογή αυτή επέτρεψε την ενοποιημένη αποθήκευση επιχειρησιακών δεδομένων, τη διαχείριση χρηστών και ρόλων, καθώς και τη δημιουργία προσαρμοσμένων \textlatin{cloud services} που συνδέουν το \textlatin{semantic middleware} επίπεδο με το τελικό περιβάλλον χρήσης. Σημαντικό μέρος αυτής της φάσης αποτελεί η χρήση \textlatin{Edge Functions}, όπως η \textlatin{\texttt{openhab-sync}}, για τη συγχρονισμένη μεταφορά και ανανέωση δεδομένων από το \textlatin{openHAB} προς το \textlatin{cloud} επίπεδο υπηρεσιών.

Η επιλογή του οικοσυστήματος \textlatin{PostgreSQL} σηματοδοτεί επίσης μια ουσιαστική μετατόπιση από την αποκλειστική λογική \textlatin{time-series persistence} προς ένα πλουσιότερο \textlatin{cloud data model}, στο οποίο τα δεδομένα του \textlatin{Digital Twin} μπορούν να συνδυάζονται με \textlatin{metadata}, ρόλους, χώρους, συσκευές και ιστορικά γεγονότα μέσα σε ένα περισσότερο \textlatin{enterprise-ready} μοντέλο.

Η πλήρης δομή της βάσης δεδομένων αποτυπώνεται στο Διάγραμμα Οντοτήτων-Σχέσεων (\textlatin{Entity-Relationship Diagram}) του Σχήματος~4.8, το οποίο περιλαμβάνει τους βασικούς πίνακες: \textlatin{organizations, organization\_members, digital\_twins, sensors, sensor\_readings, assets, events, rules, και openhab\_config}. Κάθε πίνακας φέρει \textlatin{org\_id} ως \textlatin{foreign key} για την εφαρμογή πολιτικών \textlatin{RLS}.

\subsection{\textlatin{Multi-tenancy, RBAC} και έλεγχος πρόσβασης}
Στο \textlatin{cloud} επίπεδο ενσωματώθηκε αυστηρότερη λογική ελέγχου πρόσβασης με \textlatin{Role-Based Access Control} και \textlatin{Row-Level Security}, ώστε η πλατφόρμα να υποστηρίζει απομονωμένη πρόσβαση ανά χρήστη, χώρο και οργανωτική μονάδα. Η λογική αυτή δεν σχεδιάστηκε μόνο για το παρόν \textlatin{case study}, αλλά και ως αρχιτεκτονική πρόβλεψη για πιθανή επέκταση σε περισσότερα από ένα σχολικά συγκροτήματα ή σε μελλοντικό περιβάλλον \textlatin{multi-tenant}.

Το αποτέλεσμα είναι ότι οι τεχνικές λειτουργίες διαχείρισης, οι καθημερινές προβολές διοίκησης και τα εξειδικευμένα \textlatin{panels} μπορούν να διατίθενται επιλεκτικά, με σαφή διαχωρισμό μεταξύ \textlatin{administrators}, \textlatin{operation users} και απλών χρηστών παρακολούθησης. Με αυτόν τον τρόπο, η πλατφόρμα αποκτά χαρακτήρα πραγματικής υπηρεσίας \textlatin{SaaS} και όχι απλώς απομακρυσμένου \textlatin{engineering dashboard}. Μια κεντρική επισκόπηση αυτής της νέας διεπαφής παρουσιάζεται στο Σχήμα~\ref{fig:smart_campus_dashboard}. Το Σχήμα 4.3 παρουσιάζει την κεντρική επισκόπηση (\textlatin{Overview}) της πλατφόρμας. Τέσσερα στοιχεία αξίζουν ιδιαίτερης προσοχής: πρώτον, η κάρτα κατάστασης σχολείου (Ανοιχτό/Κλειστό) που αποδεικνύει ότι το σύστημα είναι \textlatin{context-aware} και προσαρμόζεται αυτόματα στο σχολικό ωράριο· δεύτερον, ο αριθμός των \textlatin{online} αισθητήρων με \textlatin{Live} badge που επιβεβαιώνει ότι υπάρχει πραγματική \textlatin{IoT} υποδομή· τρίτον, η μηνιαία κατανάλωση ενέργειας (\textlatin{kWh}) ως κεντρικό \textlatin{KPI} που συνδέεται άμεσα με τη μείωση 19,7\% που πέτυχε το σύστημα· και τέταρτον, η χαρτογράφηση του \textlatin{campus} με \textlatin{zones} που αποδεικνύει τη χωρική επίγνωση (\textlatin{spatial awareness}) του Ψηφιακού Διδύμου.

\begin{figure}[htbp]
\centering
\includegraphics[width=\textwidth]{fig_4_2}
\caption{Κεντρική επισκόπηση της \textlatin{cloud-based} πλατφόρμας \textlatin{Smart Campus} με συγκεντρωτικούς δείκτες, ειδοποιήσεις και λειτουργικές προβολές.}
\label{fig:smart_campus_dashboard}
\end{figure}

\section{\texorpdfstring{Οπτικοποίηση και \textlatin{3D} Ψηφιακό Δίδυμο}{Οπτικοποίηση και 3D Ψηφιακό Δίδυμο}}
Η τελική μορφή της πλατφόρμας αποδίδει ιδιαίτερη σημασία στο επίπεδο οπτικοποίησης, καθώς αυτό είναι το σημείο στο οποίο το τεχνικό υπόβαθρο του \textlatin{Digital Twin} μετασχηματίζεται σε εργαλείο καθημερινής επιχειρησιακής χρήσης. Για τον λόγο αυτό, το \textlatin{visual layer} υλοποιήθηκε όχι ως παθητική συλλογή γραφημάτων, αλλά ως συνδυασμός \textlatin{2D} και \textlatin{3D} προβολών που αποτυπώνουν τον χώρο, τις συσκευές και τη ζωντανή κατάσταση της εγκατάστασης.

\subsection{\texorpdfstring{\textlatin{React Three Fiber} και \textlatin{Isometric Campus View}}{React Three Fiber και Isometric Campus View}}
Για την τρισδιάστατη απεικόνιση χρησιμοποιήθηκε το \textlatin{React Three Fiber (R3F)}, το οποίο επιτρέπει τη γεφύρωση της λογικής \textlatin{component-based} ανάπτυξης του \textlatin{React} με τη μηχανή γραφικών \textlatin{Three.js} \cite{PMNDRS2024}. Η επιλογή αυτή δικαιολογείται από τη δυνατότητα ανάπτυξης παραμετρικών τρισδιάστατων σκηνών, όπου τα κτίρια, οι χώροι και οι συσκευές δεν αποτυπώνονται απλώς ως στατικά μοντέλα, αλλά ως διαδραστικά στοιχεία που ενημερώνονται σε πραγματικό χρόνο.

Στο επίπεδο αυτό αναπτύχθηκε \textlatin{isometric campus view}, μέσω της οποίας ο χρήστης μπορεί να αποκτήσει συνολική αντίληψη της εγκατάστασης, να εντοπίσει κρίσιμα συμβάντα και να πλοηγηθεί λειτουργικά από το επίπεδο του \textlatin{campus} σε επιμέρους κτίρια και δωμάτια. Η \textlatin{3D} απεικόνιση υπηρετεί έτσι όχι μόνο αισθητικούς στόχους, αλλά κυρίως στόχους χωρικής επίγνωσης και επιχειρησιακής απόκρισης.

\subsection{\texorpdfstring{\textlatin{2D Interactive Floorplans}}{2D Interactive Floorplans}}
Παράλληλα με την τρισδιάστατη προβολή, υλοποιήθηκαν διαδραστικά \textlatin{2D floorplans} που λειτουργούν ως άμεσος μηχανισμός πλοήγησης για καθημερινά σενάρια διαχείρισης. Το \textlatin{floor plan} του σχολικού συγκροτήματος εμφανίζει τη συνολική διάταξη των χώρων και την κατάσταση κάθε αίθουσας, δηλαδή αν λειτουργεί ως ενεργό \textlatin{Digital Twin} ή όχι. Ο χρήστης μπορεί να επιλέξει οποιαδήποτε αίθουσα για να μεταβεί στην αναλυτική προβολή του αντίστοιχου υποσυστήματος.

Η συνύπαρξη \textlatin{3D} και \textlatin{2D} προβολών εξυπηρετεί διαφορετικά επίπεδα χωρικής πλοήγησης. Το Σχήμα 4.4 αποδίδει την τρισδιάστατη προβολή του \textlatin{campus}, ενώ το Σχήμα 4.5 παρουσιάζει το \textlatin{floor plan} με την κατάσταση κάθε αίθουσας --- οι χώροι που λειτουργούν ως ενεργά \textlatin{Digital Twins} επισημαίνονται διακριτά, επιτρέποντας στον χρήστη να εισέλθει σε κάθε έναν για αναλυτική παρακολούθηση.

\begin{figure}[htbp]
\centering
\includegraphics[width=\textwidth]{fig_1_1}
\caption{Τρισδιάστατη απεικόνιση του \textlatin{campus} των Εκπαιδευτηρίων Πλάτων --- ο χρήστης μπορεί να πλοηγηθεί στο μοντέλο και να επιλέξει κτίριο για αναλυτική προβολή.}
\label{fig:campus_view_3d}
\end{figure}

\begin{figure}[htbp]
\centering
\includegraphics[width=\textwidth]{Screenshot_1}
\caption{\textlatin{Floor plan} του σχολικού συγκροτήματος --- κάθε αίθουσα εμφανίζεται με την κατάστασή της (\textlatin{live Digital Twin} ή όχι) και επιτρέπει την πλοήγηση στο αντίστοιχο υποσύστημα.}
\label{fig:campus_floorplan_2d}
\end{figure}

\section{\texorpdfstring{Διαλειτουργικότητα, \textlatin{AI} και Εμπειρία Χρήστη}{Διαλειτουργικότητα, AI και Εμπειρία Χρήστη}}
Η τελική μορφή της πλατφόρμας διαφοροποιείται ουσιαστικά από ένα συμβατικό σύστημα εποπτείας, επειδή ενσωματώνει ετερογενή υποσυστήματα, επεκτείνεται με αλγοριθμικά εργαλεία και προσφέρει συνεκτική εμπειρία χρήσης σε πολλαπλές κατηγορίες χρηστών. Η ενότητα αυτή παρουσιάζει τρία από τα βασικότερα ώριμα χαρακτηριστικά της: τη διαλειτουργικότητα με \textlatin{legacy} υποδομές, την εισαγωγή λειτουργιών \textlatin{AI} και τη \textlatin{responsive} φύση της διεπαφής.

\subsection{\texorpdfstring{\textlatin{Brownfield Integration} και ενοποίηση υφιστάμενων συστημάτων}{Brownfield Integration και ενοποίηση υφιστάμενων συστημάτων}}
Η αναπτυχθείσα πλατφόρμα δεν σχεδιάστηκε με λογική αντικατάστασης όλων των προϋπαρχόντων συστημάτων. Αντίθετα, υιοθετήθηκε αρχιτεκτονική \textlatin{Brownfield Integration}, σύμφωνα με την οποία υφιστάμενες υποδομές, όπως οι αυτόματες πύλες του σχολείου, μπορούν να ενταχθούν στο ενιαίο οικοσύστημα ελέγχου καθαρά μέσω λογισμικού και διασύνδεσης. Η προσέγγιση αυτή μειώνει το κόστος μετάβασης, ενισχύει τη διαλειτουργικότητα και επιτρέπει στην πλατφόρμα να λειτουργεί ως κοινός \textlatin{operational} χώρος για πολλαπλά υποσυστήματα.

Στην ίδια λογική εντάσσεται και η δυνατότητα επέκτασης προς κάμερες, \textlatin{irrigation logic} και άλλες λειτουργίες \textlatin{campus management}. Το κεντρικό μήνυμα είναι ότι το ψηφιακό δίδυμο δεν περιορίζεται στη συλλογή αισθητηριακών δεδομένων, αλλά τείνει να εξελιχθεί σε λογισμικό ενοποίησης διαφορετικών τεχνολογικών νησίδων.

\subsection{\texorpdfstring{\textlatin{Device Management} και λειτουργικές καρτέλες ελέγχου}{Device Management και λειτουργικές καρτέλες ελέγχου}}
Η καρτέλα Συσκευών (\textlatin{Devices}) λειτουργεί ως \textlatin{hardware inventory} και \textlatin{health monitor} της \textlatin{IoT} υποδομής. Παρέχει πλήρη ορατότητα σε όλο το φυσικό \textlatin{hardware} που τροφοδοτεί το Ψηφιακό Δίδυμο, από τον \textlatin{edge server} έως τους τερματικούς αισθητήρες και τα ρελέ.

Συγκεκριμένα, η καρτέλα περιλαμβάνει: (α) κατάλογο συσκευών με κάρτα για κάθε \textlatin{device} (\textlatin{Raspberry Pi 4}, \textlatin{ESP32}, \textlatin{Shelly 1PM}, μαγνητικές επαφές, \textlatin{DFPlayer}) που εμφανίζει \textlatin{specs}, πρωτόκολλο επικοινωνίας και τοποθεσία εγκατάστασης, (β) \textlatin{live telemetry} από το \textlatin{openHAB} με πραγματικές τιμές κατάστασης (\textlatin{online/offline}) και τελευταίες μετρήσεις ανά συσκευή, (γ) 24ωρο \textlatin{uptime timeline} που αποτυπώνει οπτικά την ιστορία διαθεσιμότητας κάθε κόμβου, (δ) διάγραμμα τοπολογίας δικτύου που δείχνει τη ροή δεδομένων από τους αισθητήρες προς τον \textlatin{edge server} και το \textlatin{cloud}, και (ε) τρισδιάστατη προβολή των \textlatin{hardware components} για άμεση οπτική αναγνώριση του εξοπλισμού.

Η καρτέλα αυτή τεκμηριώνει ότι το Ψηφιακό Δίδυμο στηρίζεται σε πραγματική, επαληθεύσιμη \textlatin{IoT} υποδομή. Ο αναγνώστης μπορεί να επιβεβαιώσει ακριβώς ποιο \textlatin{hardware} χρησιμοποιήθηκε, πού τοποθετήθηκε και αν λειτουργεί σωστά σε πραγματικό χρόνο.

\begin{figure}[htbp]
\centering
\includegraphics[width=0.76\textwidth]{fig_4_5}

\vspace{0.5em}

\includegraphics[width=0.76\textwidth]{fig_4_51}
\caption{Ενδεικτική παρουσίαση των καρτελών διαχείρισης συσκευών και της διεπαφής ελέγχου πυλών στο cloud περιβάλλον της πλατφόρμας.}
\label{fig:device_management}
\end{figure}

\subsection{\texorpdfstring{Widgets ενεργειακής πρόβλεψης και \textlatin{LLM-based forecasting}}{Widgets ενεργειακής πρόβλεψης και LLM-based forecasting}}
Το σύστημα διαθέτει δύο widgets ενεργειακής πρόβλεψης. Το πρώτο (\textlatin{PredictiveEnergy}) παρέχει οπτική σύγκριση προβλεπόμενης έναντι πραγματικής κατανάλωσης βάσει ιστορικών δεδομένων με \textlatin{animated bar chart}. Το δεύτερο (\textlatin{EnergyCostPrediction}) αξιοποιεί \textlatin{LLM-based forecasting} μέσω του \textlatin{Gemini API}, το οποίο λαμβάνει τα πραγματικά \textlatin{billing data} της εγκατάστασης και παράγει δομημένες προβλέψεις (\textlatin{JSON}) για ορίζοντα 3 μηνών, λαμβάνοντας υπόψη το ελληνικό τιμολόγιο ενέργειας (0,162 €/\textlatin{kWh}), το σχολικό ημερολόγιο, την εποχικότητα και τον επαληθευμένο συντελεστή μείωσης 19,7\%.

Η προσέγγιση αυτή αποτελεί \textlatin{LLM-as-a-forecaster} με \textlatin{domain knowledge}, όχι κλασικό στατιστικό μοντέλο τύπου \textlatin{ARIMA} ή \textlatin{LSTM}. Η λειτουργία αυτή οπτικοποιείται στο Σχήμα~\ref{fig:ai_energy_forecasting}.

\begin{figure}[htbp]
\centering
\includegraphics[width=0.78\textwidth]{fig_4_6}

\vspace{0.5em}

\includegraphics[width=0.78\textwidth]{fig_4_61}
\caption{Ενδεικτική απεικόνιση των δύο widgets ενεργειακής πρόβλεψης της πλατφόρμας: \textlatin{PredictiveEnergy} και \textlatin{EnergyCostPrediction}.}
\label{fig:ai_energy_forecasting}
\end{figure}

\subsection{\texorpdfstring{\textlatin{Responsive} πρόσβαση και χρήση από μη τεχνικούς ρόλους}{Responsive πρόσβαση και χρήση από μη τεχνικούς ρόλους}}
Σε αντίθεση με την πρώτη φάση, κατά την οποία η απομακρυσμένη τεχνική πρόσβαση στο τοπικό πρωτότυπο απαιτούσε περισσότερο εξειδικευμένη δικτυακή διαχείριση, η τελική \textlatin{cloud} πλατφόρμα σχεδιάστηκε ώστε να είναι άμεσα προσβάσιμη μέσω \textlatin{browser}, χωρίς να απαιτείται \textlatin{VPN} για τον τελικό χρήστη. Η \textlatin{responsive} φύση της εφαρμογής επιτρέπει χρήση από κινητά τηλέφωνα, \textlatin{tablets} και σταθμούς εργασίας, διασφαλίζοντας ότι η πλατφόρμα μπορεί να χρησιμοποιηθεί καθημερινά από διευθυντές και προσωπικό λειτουργίας, όχι μόνο από μηχανικούς ή τεχνικούς.

Αυτό το χαρακτηριστικό δεν αποτελεί απλή βελτίωση εργονομίας. Αντιπροσωπεύει μία ουσιώδη αλλαγή στο ίδιο το \textlatin{operational} μοντέλο της πλατφόρμας. Η πληροφορία και ο έλεγχος γίνονται πλέον διαθέσιμα στο σημείο λήψης αποφάσεων, γεγονός που αυξάνει τη χρηστικότητα, επιταχύνει την απόκριση σε συμβάντα και διευρύνει σημαντικά το πεδίο πρακτικής υιοθέτησης της πλατφόρμας. Η \textlatin{responsive} συμπεριφορά της πλατφόρμας απεικονίζεται στο Σχήμα~\ref{fig:responsive_views}.

\begin{figure}[htbp]
\centering
\includegraphics[width=0.31\textwidth]{fig_4_4}\hfill
\includegraphics[width=0.31\textwidth]{fig_4_41}\hfill
\includegraphics[width=0.31\textwidth]{fig_4_42}
\caption{Απεικόνιση της \textlatin{responsive} συμπεριφοράς της πλατφόρμας σε κινητές και φορητές συσκευές.}
\label{fig:responsive_views}
\end{figure}

\section{Συνολική Τεχνική Αποτίμηση του Κεφαλαίου}
Το παρόν κεφάλαιο περιέγραψε τη μετάβαση από το τοπικό πρωτότυπο με \textlatin{MQTT}, \textlatin{openHAB}, \textlatin{InfluxDB} και \textlatin{Grafana} σε μια πολυεπίπεδη αρχιτεκτονική \textlatin{cloud SaaS} με \textlatin{React}, \textlatin{Supabase/PostgreSQL} και \textlatin{Edge Functions}. Η τελική υλοποίηση συνδυάζει \textlatin{edge} συλλογή δεδομένων, σημασιολογική οργάνωση, αμφίδρομο έλεγχο, χωρική οπτικοποίηση \textlatin{2D/3D}, διαβαθμισμένη πρόσβαση και εξειδικευμένα widgets λειτουργικής ανάλυσης. Από αρχιτεκτονική άποψη, το σύστημα διαχωρίζει το επίπεδο πεδίου από το επίπεδο υπηρεσιών και διεπαφής, ώστε να υποστηρίζεται κλιμάκωση, απομακρυσμένη πρόσβαση και επέκταση σε πρόσθετα υποσυστήματα του σχολικού συγκροτήματος.

Η τελική μορφή της πλατφόρμας είναι επιχειρησιακά ενεργή στη διεύθυνση \href{https://digitaltwin.pi-tech.gr}{\textlatin{digitaltwin.pi-tech.gr}}. Η πρόσβαση στα πραγματικά δεδομένα της εγκατάστασης παρέχεται αποκλειστικά σε εξουσιοδοτημένους χρήστες μέσω διαπιστευτηρίων. Παράλληλα, διατίθεται \textlatin{Demo} περιβάλλον με προσομοιωμένα δεδομένα και πλήρη διαδραστικότητα --- συμπεριλαμβανομένου του ελέγχου πυλών και της \textlatin{quiet-lamp} --- το οποίο είναι προσβάσιμο μέσω του κουμπιού \textlatin{'Demo Sign In'} στη σελίδα εισόδου, χωρίς να απαιτούνται διαπιστευτήρια. Το \textlatin{Demo} περιβάλλον επιτρέπει την αξιολόγηση της πλατφόρμας και την κατανόηση του κύκλου \textlatin{Sense--Decide--Act} χωρίς πρόσβαση σε πραγματικά δεδομένα ή συσκευές.

